\documentclass[12pt]{article}
\usepackage[T1]{fontenc}
\usepackage[utf8]{inputenc}
\usepackage{lmodern}
\usepackage{textcomp}
\usepackage{enumitem}
\usepackage{geometry}
\geometry{margin=1in}
\begin{document}
\begin{center}
Answer Key - Political Science - Undergraduate Level Assessment 24 \
\small (Answers are ordered exactly as in the question paper.)
\end{center}

\section*{Multiple Choice}
\begin{enumerate}[label=\arabic*.]
\item \textbf{Answer:} A
\\ \textit{Options:} A) contribute money to candidates for election; B) coordinate local get-out-the-vote campaigns; C) promote the defeat of incumbents in the federal and state legislatures; D) organize protest demonstrations and other acts of civil disobedience
\\ \textit{Note:} The primary function of political action committees (PACs) is to influence the electoral process by raising and contributing money to candidates, thereby enabling them to fund their campaigns and gain political office.
\item \textbf{Answer:} D
\\ \textit{Options:} A) Ways and Means; B) Judiciary; C) Ethics; D) Rules
\\ \textit{Note:} The Rules Committee in the House of Representatives is responsible for establishing the procedures for how bills are debated and amended, making it the key committee that governs legislative process.
\item \textbf{Answer:} A
\\ \textit{Options:} A) A federal mandate; B) A constitutional amendment; C) Affirmative action; D) Tort reform
\\ \textit{Note:} The No Child Left Behind Act is a federal mandate because it requires states to adhere to specific educational standards in order to qualify for federal funding, thereby imposing a requirement that states must follow to receive financial support.
\item \textbf{Answer:} B
\\ \textit{Options:} A) It encouraged greater involvement in European politics; B) It was designed as the antithesis of European politics; C) It created a large standing army; D) It encouraged the centralization of political power in the US
\\ \textit{Note:} The 'Philadelphian System' was conceived to reflect American ideals of democracy and individual rights, positioning itself as a distinct alternative to the monarchies and aristocracies prevalent in Europe, thereby embodying the notion of American exceptionalism.
\item \textbf{Answer:} A
\\ \textit{Options:} A) Withdrawal of economic trade rights with the domestic market.; B) Export controls protecting technological advantage and further foreign policy objectives.; C) Control of munitions and arms sales.; D) Import restrictions to protect a domestic market from foreign goods.
\\ \textit{Note:} The withdrawal of economic trade rights with the domestic market is the odd one out because it directly impacts domestic economic relations rather than serving as a tool for influencing international security policy, unlike the other options which typically focus on external economic sanctions or incentives.
\end{enumerate}

\section*{True / False}
\begin{enumerate}[label=\arabic*.]
\setcounter{enumi}{5}
\item \textbf{Answer:} False
\\ \textit{Options:} A) True; B) False
\\ \textit{Note:} The salience of a political issue is influenced by various factors including public perception, social movements, and political agendas, rather than being solely determined by media coverage.
\item \textbf{Answer:} False
\\ \textit{Options:} A) True; B) False
\\ \textit{Note:} The statement is false because the responsibility to prevent is a component of the broader responsibility to protect doctrine, which includes three pillars: the responsibility to prevent, to react, and to rebuild, indicating that prevention alone does not encompass the entire doctrine.
\item \textbf{Answer:} True
\\ \textit{Options:} A) True; B) False
\\ \textit{Note:} Autocracy, which concentrates power in a single authority and undermines individual freedoms, contradicts the fundamental principles of the 'American Creed' that prioritize liberty, equality, and individualism.
\item \textbf{Answer:} True
\\ \textit{Options:} A) True; B) False
\\ \textit{Note:} The phrase 'empire by invitation' describes a situation where a state willingly seeks the assistance or protection of a more powerful external entity, thereby establishing a form of dependency that is characterized by voluntary reliance for security rather than coercion.
\item \textbf{Answer:} True
\\ \textit{Options:} A) True; B) False
\\ \textit{Note:} Since 1972, the rise of televised debates, social media, and personalized campaign strategies has led voters to prioritize individual candidates' personalities and appeal over traditional party affiliations.
\end{enumerate}

\section*{Long Form Response}
\begin{enumerate}[label=\arabic*.]
\setcounter{enumi}{10}
\item \textbf{Answer:} Offshore Balancing is a grand strategy that emphasizes the importance of embracing multi- polarity and encouraging greater restraint in international affairs. This approach suggests that the United States and other major powers should avoid direct military interventions and instead allow regional powers to take on greater responsibility for their own security. By doing so, offshore balancing aims to reduce the burden on the intervening state while still maintaining stability in critical regions. The implications for international relations are significant, as this strategy can lead to a more balanced distribution of power and encourage states to enhance their security capabilities, ultimately fostering a more self-reliant global order. Such a shift also necessitates a reevaluation of alliances and partnerships, as states must become more accountable for their security, thereby reshaping the dynamics of international cooperation and conflict.
\\ \textit{Note:} correct because it accurately outlines offshore balancing as a strategy that promotes multi-polarity and restraint in military interventions, while emphasizing regional powers' roles in maintaining their own security, thereby reducing the intervening state's burden.
\item \textbf{Answer:} American political parties, primarily the Democratic and Republican parties, often exhibit less ideological consistency compared to parties in multiparty systems like England and Israel. In the United States, party identification can be fluid, with individuals aligning with parties based on specific issues rather than a cohesive ideological framework. For example, moderate Democrats and conservative Republicans may hold divergent views on key topics, reflecting a broader spectrum of beliefs within each party. In contrast, multiparty systems tend to have parties that are more clearly defined by consistent ideological positions, allowing voters to make more distinct choices based on their political beliefs. This ideological variability in the U.S. leads to a more complex and sometimes fragmented political landscape, where party affiliation does not always equate to a uniform set of principles.
\\ \textit{Note:} correct because it highlights the fluidity of party identification in the American political system, which allows for a wider range of ideological diversity within the two dominant parties, contrasting with the more ideologically consistent and specialized parties typically found in multiparty systems like England and Israel.
\end{enumerate}

\vfill
\noindent\textit{End of Answer Key}
\end{document}